Atkārtotajā CPU eksperimentā (Intel Core i7-1355U) vērtēti pirmie 200 validācijas piemēri ar palīgpunktu. Stingrajā salīdzinājumā saglabāts jaunā punkta vārds un predikātu argumentu secība. Predikātu savstarpējā secība rezultātu neietekmē. No 200 piemēriem 13 atsauces bija bez predikātiem. Tās ieskaitītas kā nesakritības. Modelim \texttt{mid} pirmais piedāvājums sakrita 99 gadījumos, vismaz viens no četriem -- 129 gadījumos. Iepriekš lietotais brīvākais kritērijs tajās pašās prognozēs deva 136 sakritības no 200. Pirmajos 50 piemēros stingrais četru piedāvājumu rādītājs bija 32/50 modelim \texttt{mid}, 33/50 modelim GenesisGeo un 4/50 modelim SmolLM2-360M. Atkārtotajā JGEX pārbaudē DD-AR atrisināja 205 uzdevumus, bet kopā ar atsevišķo palīgpunktu meklējumu rezultātiem pierādīti 217 no 231. IMO-AG-30 pieejamajā daļā ar dedukciju pierādīti 15, bet pēc papildu meklējumiem -- 16 no 29 uzdevumiem.

Modeļa \texttt{small} 6300. apmācības solī pirmā prognoze sakrita ar etalonu 92 no 200 piemēriem. Vērtējot visus četrus piedāvājumus, sakritību skaits pieauga līdz 137. Uzdevumos, kurus sākotnējā dedukcija neatrisināja, ar šo modeli pierādīti 6 no 26 JGEX un 1 no 14 IMO uzdevumiem. Kontrolpunktam ar zemāko validācijas zudumu (20000 soļi) pirmā prognoze sakrita ar etalonu 113 no 200 gadījumiem, bet vismaz viena no četrām -- 143. Ar to pierādīti 6 no 26 sākotnēji neatrisinātajiem JGEX un 0 no 14 IMO uzdevumiem. Atsevišķajā 101 uzdevuma palīgkonstrukciju kopā dedukcija vien pierādīja 0 uzdevumus, bet trīs papildu meklējumu rezultātu apvienojums -- 11. Piemēram, 2004. gada atlases kopas uzdevumā G4a heiristika pievienoja trijstūrim $ABC$ apvilktās riņķa līnijas centru $E$. No rādiusu vienādības $AE=CE$ un leņķu sakarībām iegūst $\angle AEP=\angle PEC$. Trijstūriem ir kopīga mala $EP$. Trijstūri $AEP$ un $CEP$ ir vienādi pēc divām malām un leņķa starp tām, tātad $AP=CP$. AlphaGeometry lokāli pabeidza trīs pārbaudes piemērus (sakritība četros piedāvājumos 1/3). Šī izlase ir par mazu modeļu salīdzināšanai. Katrā palīgpunkta prognozes uzdevumā bija dota viena etalonkonstrukcija. Iespējami arī citi derīgi risinājumi.

Starptautiskās matemātikas olimpiādes dalībnieku rezultātus ar dedukciju salīdzināja 6 gados, kuros vietējā kopā bija vismaz divi dažādi tās pašas olimpiādes uzdevumi (kopā 12). Vidējo vērtējumu aprēķināja visiem dalībniekiem, ieskaitot nulles. Trijos pāros dedukcija pierādīja vienu no abiem uzdevumiem. Visos 3 gadījumos šim uzdevumam dalībnieku vidējais vērtējums bija augstāks. Pārējos 3 pāros abu uzdevumu iznākums sakrita. Piemēram, 2015. gada 3. un 4.~uzdevuma vidējie vērtējumi bija 0,65 un 4,79 no septiņiem punktiem, un dedukcija pierādīja otro no šīm formalizācijām. \href{https://www.imo-official.org/results/individual/year/2015/}{(Starptautiskā matemātikas olimpiāde, 2015)}.